%% kapa: exam=1 round=29 year=2018 period=1 form=A subject=economics subject_ko=경제학원론 seq=7 number=47
%% source: Q-Net 감정평가사 기출문제 (cst00309) artlSeq=5202808 file=extracted/1008/5202808/2018년 제29회 감정평가사 1차시험 기출문제/2018년 제29회 감정평가사 1차시험 기출문제/1교시 A형.pdf
%% conversion: kordoc 4.24.51 -> tools/gen1.py (auto); review=proofread-sonnet; flags=none
\begin{question}{47}
가격차별의 사례가 아닌 것은?
\begin{choices}
  \item 영화관 일반 요금은 1만 원, 심야 요금은 8천 원이다.
  \item 놀이공원 입장료는 성인 5만 원, 청소년 3만 원이다.
  \item 동일한 롱패딩 가격은 겨울에 30만 원, 여름에 20만 원이다.
  \item 동일한 승용차 가격은 서울에서 2,000만 원, 제주에서 1,500만 원이다.
  \item 주간 근무자 수당은 1만 원, 야간 근무자의 수당은 1만5천 원이다.
\end{choices}
\end{question}
\answer{5}
